
\subsection{博特--莫尔斯情形下相容多值截面系统的归纳构造}
\label{subsec:7-2-2}

本小节先用仓西结构、其纤维积及多值截面的语言，
重新表述上一小节的问题，
再给出这一问题的详细解法。
先考虑如下情形。
设 $(\mathfrak B,\le)$ 为偏序集，
且对每个 $\beta\in\mathfrak B$ 满足有限性条件
\[
\#\{\beta'\in\mathfrak B\mid\beta'\le\beta\}<\infty.
\]

\begin{statement}{情形}{7.2.2}\label{sit:7-2-2}
对每个 $\beta\in\mathfrak B$，
有一个带仓西结构的空间 $X(\beta)$，满足：

\noindent\textup{(7.2.3.1)}\quad
$X(\beta)$ 有定向切丛。

\noindent\textup{(7.2.3.2)}\quad
$X(\beta)$ 的余维 $m$ 角
$S_m(X(\beta))$
作为带仓西结构的空间，
同构于下列直积的并：
\[
S_m(X(\beta))
=\bigcup_{i=1}^{N_m(\beta)}
\prod_{j=0}^{m}X(\beta_{i,j}),
\qquad \beta_{i,j}<\beta.
\]
\end{statement}

在情形~\ref{sit:7-2-2} 下可证明：
\begin{lemma}{7.2.4}\label{lem:7-2-4}
存在 $X(\beta)$ 上的一族多值截面
$\{s_\beta\}_{\beta\in\mathfrak B}$，
满足：

\noindent\textup{(7.2.5.1)}\quad
在
$\prod_{j=0}^{m}X(\beta_{i,j})
\subset S_m(X(\beta))$
这一层上，$s_\beta$ 的限制等于
$\bigoplus_{j=0}^{m}s_{\beta_{i,j}}$。

\noindent\textup{(7.2.5.2)}\quad
$s_\beta$ 与零截面横截。
\end{lemma}

证明由定理 A1.23
（\textbf{[FuOn99II]} 引理 3.14）
和下列显然的引理立即得到。

\begin{lemma}{7.2.6}\label{lem:7-2-6}
设各 $X_j$ 为带仓西结构的空间，
各 $s_j$ 为其上的多值截面。
若每个 $s_j$ 都与零截面横截，
则在 $\prod_jX_j$ 上，
$\bigoplus_js_j$ 也与零截面横截。
\end{lemma}

引理~\ref{lem:7-2-4} 的证明本身
只用到了 \textbf{[FuOn99II]} 第 1 章的内容。
但我们实际面对的几何情形
并非情形~\ref{sit:7-2-2}，
而是下面的情形：
必须考虑的是\emph{纤维积}，
而不是直积。

\begin{statement}{情形}{7.2.7}\label{sit:7-2-7}
对每个 $\beta\in\mathfrak B$，
有一个带仓西结构的空间 $X(\beta)$，
以及到光滑流形 $L(\beta)$ 的
弱浸没且强连续的映射
$\operatorname{ev}_\beta\colon
X(\beta)\to L(\beta)$。
$L(\beta)$ 上还有实线丛
$\Theta_{L(\beta)}$。
它们具有下列性质：

\noindent\textup{(7.2.8.1)}\quad
$X(\beta)$ 有切丛，且
\[
\det TX(\beta)\otimes
\operatorname{ev}_\beta^*
\bigl(\Theta_{L(\beta)}\otimes
\det TL(\beta)\bigr)
\]
带有一个平凡化。

\noindent\textup{(7.2.8.2)}\quad
$S_m(X(\beta))$ 是若干
$X(\beta')$（$\beta'<\beta$）
的纤维积之并。
更准确地说，
存在
$\beta_{i,j}<\beta$
（$i=1,\ldots,N(\beta)$，
$j=0,\ldots,m$）
及直积分解
\begin{equation}
L(\beta_{i,j})
=L_1(\beta_{i,j})\times L_2(\beta_{i,j}),
\tag{7.2.9}\label{eq:7-2-9}
\end{equation}
使
\[
S_mX(\beta)
=\bigcup_{i=1}^{N(\beta)}
\left(
\prod_{j=0}^{m}X(\beta_{i,j})
{}_{\operatorname{ev}}\!\times_{\prod_{j=0}^{m}L_1(\beta_{i,j})}
L'(\beta,m,i)
\right).
\]
这里
$L'(\beta,m,i)\subset
\prod_{j=0}^{m}L_1(\beta_{i,j})$
是子流形，
而映射
\[
\operatorname{ev}\colon
\prod_{j=0}^{m}X(\beta_{i,j})
\longrightarrow
\prod_{j=0}^{m}L_1(\beta_{i,j})
\]
由各因子的
$\operatorname{ev}_{\beta_{i,j}}$
以显然的方式诱导。

\noindent\textup{(7.2.8.3)}\quad
有同构
\[
L(\beta)\cong
\prod_{j=0}^{m}L_2(\beta_{i,j}).
\]
在纤维积层
\[
\prod_{j=0}^{m}X(\beta_{i,j})
{}_{\operatorname{ev}}\!\times_{\prod_jL_1(\beta_{i,j})}
L'(\beta,m,i)
\subset S_mX(\beta)
\]
上，
$\operatorname{ev}_\beta\colon
X(\beta)\to L(\beta)$
的限制，
正是由各因子的
$\operatorname{ev}_{\beta_{i,j}}$
诱导的到
$\prod_{j=0}^{m}L_2(\beta_{i,j})$
的映射。

\noindent\textup{(7.2.8.4)}\quad
条件 (7.2.8.1) 的平凡化
与 (7.2.8.2)、(7.2.8.3)
的同构相容。
\end{statement}

命题~\ref{prop:7-1-1}、
\ref{prop:7-1-2} 表明，
取
$X((\beta,k))
=\mathcal M^{\mathrm{main}}_{k+1}(\beta)$
时便出现情形~\ref{sit:7-2-7}。
这里 $\mathfrak B$ 是对子 $(\beta,k)$
按字典序排列所得的集合，
且 $L((\beta,k))=L^{k+1}$。

本节构造多值截面的主要困难在于，
上述针对\emph{直积}的
引理~\ref{lem:7-2-4} 的证明
无法推广到情形~\ref{sit:7-2-7}
中的\emph{纤维积}。
这正是上一小节所说的困难。
例如，对于格罗莫夫--威腾不变量，
无须担心类似问题，
因为那里只须分别为每个 $\beta$
构造 $X(\beta)$ 的虚基本闭链，
无需顾及不同 $\beta$ 的
$X(\beta)$ 上扰动的相容性。
综上，困难源于我们必须在\emph{链层面}
实施虚构造；这是主要障碍。

为克服此障碍，
我们将仔细选取 $L$ 的奇异单形集，
使扰动依赖于 $L$ 的奇异单形
（或情形~\ref{sit:7-2-7} 中
$L_1(\beta)$ 的奇异单形）。

先回顾 \textbf{[Fuk96II]}
中的博特--莫尔斯理论模型\liftIndexRangeStart{I384}，
那里已出现这样的构造。
设 $M$ 是有限维紧致、定向、
无边界的黎曼流形，
$f\colon M\to\mathbb R$
是博特--莫尔斯函数。
将 $f$ 的临界点集
$\operatorname{Cr}(f)$
分解为连通分支
\[
\operatorname{Cr}(f)
=\bigcup_{i=1,2,\ldots}R_i.
\]
按 $i<j$ 蕴含
$f(R_i)\le f(R_j)$ 的方式编号。
从 $R_i$ 到 $R_j$ 的连接轨线集为
\[
\widetilde{\mathcal M}(R_i,R_j)
=\left\{
\ell\colon\mathbb R\to M\ \middle|\
\frac{d\ell}{dt}=\operatorname{grad}f,
\ \lim_{t\to-\infty}\ell(t)\in R_i,
\ \lim_{t\to+\infty}\ell(t)\in R_j
\right\}.
\]
它有由
$\ell(\cdot)\mapsto\ell(t+\cdot)$
给出的 $\mathbb R$ 作用。
记商空间为
$\mathring{\mathcal M}(R_i,R_j)$。
容易看出，
可在临界点集之外扰动给定的
博特--莫尔斯函数 $f$，
使 $\mathring{\mathcal M}(R_i,R_j)$
成为光滑流形。

为简便，
假设 $R_i$ 及其上的稳定流形丛均已定向。
一般情形可仿照第~3.7.5 小节，
在 $R_i$ 上取局部系数系。
以 $\omega$-极限点与 $\alpha$-极限点定义
\[
\operatorname{ev}^{i,j}
=\bigl(\operatorname{ev}^{i,j}_{-\infty},
\operatorname{ev}^{i,j}_{+\infty}\bigr)
\colon\mathring{\mathcal M}(R_i,R_j)
\longrightarrow R_i\times R_j,
\qquad
\operatorname{ev}^{i,j}_{\pm\infty}[\ell]
=\lim_{t\to\pm\infty}\ell(t).
\]
再在临界点集之外扰动 $f$，
可假设
$\mathring{\mathcal M}(R_i,R_j)$
有紧化 $\mathcal M(R_i,R_j)$，
其余维 $m$ 的角为
\begin{equation}
\begin{aligned}
S_m\mathcal M(R_i,R_j)
={}&\bigcup_{i=\ell_1<\cdots<\ell_{m+2}=j}\\
&\mathring{\mathcal M}(R_{\ell_1},R_{\ell_2})
\mathop{\times}_{R_{\ell_2}}\cdots
\mathop{\times}_{R_{\ell_{m+1}}}
\mathring{\mathcal M}(R_{\ell_{m+1}},R_{\ell_{m+2}}).
\end{aligned}
\tag{7.2.10}\label{eq:7-2-10}
\end{equation}
于是落在情形~\ref{sit:7-2-7} 中。
还要注意：
\begin{equation}
i>j\quad\Longrightarrow\quad
\mathcal M(R_i,R_j)=\varnothing.
\tag{7.2.11}\label{eq:7-2-11}
\end{equation}

约定 $i=j$ 时
$\mathcal M(R_i,R_i)=\varnothing$。
利用~\eqref{eq:7-2-11}
可以证明下面的引理。
记 $S(R_i;\mathbb Z)$ 为
$R_i$ 上的光滑奇异链复形。

\begin{lemma}{7.2.12}\label{lem:7-2-12}
（参见 \textbf{[Fuk96II]} 第 97 页）
存在 $R_i$ 上光滑奇异单形的
可数集 $\mathcal X(R_i)$\liftIndex{I239}，
它生成 $S(R_i;\mathbb Z)$
的子复形 $C(R_i;\mathbb Z)$，
满足下列性质：

\noindent\textup{(7.2.13.1)}\quad
自然包含
$C(R_i;\mathbb Z)\to S(R_i;\mathbb Z)$
是链同伦等价。

\noindent\textup{(7.2.13.2)}\quad
对 $\mathcal X(R_i)$ 的每个元素
$(P,f)$（$P$ 为单形，
$f\colon P\to R_i$ 为映射），
$f$ 均与
\[
\operatorname{ev}^{i,j}_{-\infty}
\colon S_m\mathcal M(R_i,R_j)
\longrightarrow R_i
\]
横截；这里 $j$ 任意，
$m=0,1,2,\ldots$。

\noindent\textup{(7.2.13.3)}\quad
对 $\mathcal X(R_i)$ 的每个元素
$(P,f)$，纤维积所得的链
\[
\bigl(
P_f\mathop{\times}_{\operatorname{ev}^{i,j}_{-\infty}}
S_m\mathcal M(R_i,R_j),
\operatorname{ev}^{i,j}_{+\infty}
\bigr)
\in C(R_j;\mathbb Z)
\]
都可分解为 $\mathcal X(R_j)$
中的元素；这里 $j$ 任意，
$m=0,1,2,\ldots$。
\translatornote{原书此处把输出单形集写成
$\mathcal X(R_i)$，
但同一条件已写明输出属于
$C(R_j;\mathbb Z)$，
且输出评价映射的靶为 $R_j$；
这里按输出型别写 $\mathcal X(R_j)$。
参见疑误 A-0432。}
\end{lemma}

这里“$f\colon P\to R_i$
与 $\operatorname{ev}^{i,j}_{-\infty}$ 横截”
意指：对 $P$ 的每个任意余维面
$P_a$，限制
$f|_{P_a}\colon P_a\to R_i$
均与该评价映射横截。

\noindent\textit{证明。}
对 $i$ 归纳构造 $\mathcal X(R_i)$。
标准横截性定理立即给出满足
(7.2.13.1)、(7.2.13.2) 的
$\mathcal X(R_1)$。
设已经对 $i\le i_0$
选定 $\mathcal X(R_i)$，
使 (7.2.13.1) 对每个
$i\le i_0$ 成立，
(7.2.13.2) 对这些 $i$ 及任意 $j$ 成立，
(7.2.13.3) 对 $i,j\le i_0$ 成立。
现在构造 $\mathcal X(R_{i_0+1})$。
先考虑纤维积链
\begin{equation}
\bigl(
P_f\mathop{\times}_{\operatorname{ev}^{i,i_0+1}_{-\infty}}
S_m\mathcal M(R_i,R_{i_0+1}),
\operatorname{ev}^{i,i_0+1}_{+\infty}
\bigr).
\tag{7.2.14}\label{eq:7-2-14}
\end{equation}
由归纳假设及~\eqref{eq:7-2-11}，
式~\eqref{eq:7-2-14}
中的纤维积是横截的。
还需注意，
式~\eqref{eq:7-2-14} 的输出
与
\[
\operatorname{ev}^{i_0+1,j}_{-\infty}
\colon S_{m'}\mathcal M(R_{i_0+1},R_j)
\longrightarrow R_{i_0+1}
\]
对任意 $j$ 及
$m'=0,1,2,\ldots$ 均横截。
后者由归纳假设推出：
\[
S_m\mathcal M(R_i,R_{i_0+1})
\mathop{\times}_{R_{i_0+1}}
S_{m'}\mathcal M(R_{i_0+1},R_j)
\]
就是
$S_{m+m'+1}\mathcal M(R_i,R_j)$
的一个分量。
对式~\eqref{eq:7-2-14} 的空间
作单纯剖分，
使所得每个单形仍满足
(7.2.13.2) 的逐面横截性。
\translatornote{原书此处写“仍满足 (7.2.11)”，
但 (7.2.11) 是模空间为空的指标条件，
不是单形性质；
这里按紧前的横截性论证改指 (7.2.13.2)。
参见疑误 A-0431。}
该剖分可与边界上先前所取的剖分相容。
各剖分单形都可视作奇异单形；
记由此得到的全部单形之集为
$\mathcal X_0(R_{i_0+1})$。

以下选取包含
$\mathcal X_0(R_{i_0+1})$
的 $\mathcal X(R_{i_0+1})$。
设 $S_{\mathrm{trans}}(R_{i_0+1};\mathbb Z)$
为由集合
$\mathcal X_{\mathrm{trans}}(R_{i_0+1})$
生成的子复形；
后一个集合包含全部奇异单形
$(P,f)$，
它们及其任意余维的所有面
均满足 (7.2.13.2)、(7.2.13.3)。
原书称由标准横截性论证可得，
包含映射
\[
S_{\mathrm{trans}}(R_{i_0+1};\mathbb Z)
\longrightarrow S(R_{i_0+1};\mathbb Z)
\]
是链同伦等价。
\translatornote{这里保留原书同时要求
(7.2.13.2)、(7.2.13.3) 的读法。
后一条件涉及尚未选定的输出单形集
$\mathcal X(R_j)$；
仅凭所称标准横截性论证，
对这个附加闭包条件的链同伦等价
尚未得到证明。
参见疑误 A-0434。}
然后利用标准代数论证，
如下选出 $\mathcal X(R_{i_0+1})$：
先取
$\mathcal X_1(R_{i_0+1})$
使
\[
\mathcal X_0(R_{i_0+1})
\subset\mathcal X_1(R_{i_0+1})
\subset\mathcal X_{\mathrm{trans}}(R_{i_0+1}),
\]
由它生成子复形
$C_1(R_{i_0+1})$，
并使
$H(C_1(R_{i_0+1}))\to H(R_{i_0+1})$
满射。
若此映射不是单射，
就取
$\mathcal X_2(R_{i_0+1})
\supset\mathcal X_1(R_{i_0+1})$，
仍包含于
$\mathcal X_{\mathrm{trans}}(R_{i_0+1})$，
使它生成子复形
$C_2(R_{i_0+1})$，
并使
\[
\ker\bigl(
H(C_1(R_{i_0+1}))
\longrightarrow H(R_{i_0+1})
\bigr)
\]
在 $H(C_2(R_{i_0+1}))$ 中变成零。
若
$H(C_2(R_{i_0+1}))\to H(R_{i_0+1})$
仍有非零核，
便同样取
$\mathcal X_3(R_{i_0+1})$，
如此继续。
于是
\[
\mathcal X(R_{i_0+1})
=\bigcup_k\mathcal X_k(R_{i_0+1})
\]
具有所需性质。
引理~\ref{lem:7-2-12} 证毕。

\begin{remark}{7.2.15}\label{rem:7-2-15}
\textup{(1)}
式~\eqref{eq:7-2-11}
在上面的证明中起关键作用：
取定 $\mathcal X(R_i)$ 以后，
构造 $j>i$ 的
$\mathcal X(R_j)$ 时
不必再修改它。

\textup{(2)}
如果不要求 $C(R_i;\mathbb Z)$
由奇异单形生成，
便可以取它为有限生成子复形。

\textup{(3)}
即使给定单纯剖分中一个单形的定向
及仿射结构，
它与标准单形的识别仍不唯一。
为排除自同构，
须指定顶点的次序。
这也是 S. Eilenberg
创立奇异同调论时的重要问题之一
（见 \textbf{[Eil44]} 导言）。
可以如下固定顶点次序。

待剖分的空间是带角流形。
按层的维数归纳固定顶点次序，
并要求：
若 $v_1$ 位于 $d_1$ 维层内部，
$v_2$ 位于 $d_2$ 维层内部，
且 $d_1<d_2$，
则 $v_1<v_2$。
依归纳假设，
同一层中的顶点次序已定；
对边界层这是归纳假设，
对内部顶点则可任意固定其次序。
再利用字典序，
即可比较两个顶点，
除非它们处于不同但同维的层。
将同维各层的次序也固定，
余下情形的顶点次序便随之确定。

事实上，
只要与归纳较早阶段
在各边界层已固定的次序相容，
其余次序可以自由选择。
固定顶点次序后，
便可据此将单形识别为标准单形，
再按此识别是否保持定向赋予符号。
这样就为模空间与各 $P$
的纤维积定义基本链。
\end{remark}

\textbf{[Fuk96II]}
利用类似引理~\ref{lem:7-2-12}
的结果构造了谱序列
\begin{equation}
E^2_{ij}=H_i(R_j;\mathbb Z)
\Longrightarrow H(M;\mathbb Z).
\tag{7.2.16}\label{eq:7-2-16}
\end{equation}
其方法如下
（见 \textbf{[Fuk96II]} 定理 1.2）。
不过，谱序列~\eqref{eq:7-2-16}
本身也直接由 \textbf{[Bot59]} 得到；
这里的新意是下面的论证
适于推广到无限维情形，
特别是弗洛尔同调。
令
\[
C(M,f)=\bigoplus_i C(R_i;\mathbb Z),
\]
并定义
\begin{equation}
\partial^+(P,f)
=\partial(P,f)
+\sum_{j\ge i}
\bigl(
P_f\mathop{\times}_{\operatorname{ev}^{i,j}_{-\infty}}
\mathcal M(R_i,R_j),
\operatorname{ev}^{i,j}_{+\infty}
\bigr),
\tag{7.2.17}\label{eq:7-2-17}
\end{equation}
其中 $(P,f)\in\mathcal X(R_i)$，
$\partial\colon C(R_i;\mathbb Z)
\to C(R_i;\mathbb Z)$
是通常的边界算子。
这里暂不讨论符号。
可以证明
$\partial^+\circ\partial^+=0$，
且
$(C(M,f),\partial^+)$
的同调群同构于 $M$ 的同调群。
再用滤过
\[
F^I C(M,f)
=\bigoplus_{i\ge I}C(R_i;\mathbb Z)
\]
便得到谱序列~\eqref{eq:7-2-16}。

独立于 \textbf{[Fuk96II]}，
D. Austin 与 P. Braam
在 \textbf{[AuBr95]}
研究了一个在 $\mathbb R$ 系数下
类似~\eqref{eq:7-2-16}
的谱序列，
但须加上下列假设。

\begin{statement}{条件}{7.2.18}\label{cond:7-2-18}
（\textbf{[AuBr95]} 第 138 页）
映射
\[
\operatorname{ev}^{i,j}\colon
\mathring{\mathcal M}(R_i,R_j)
\longrightarrow R_i\times R_j
\]
是局部平凡纤维丛。
\end{statement}

应注意，
对给定的博特--莫尔斯函数 $f$，
在临界点集外对 $f$ 作一般扰动，
通常并不满足条件~\ref{cond:7-2-18}。
因此，\textbf{[AuBr95]} 的结果
不足以为一般的博特--莫尔斯函数
构造谱序列~\eqref{eq:7-2-16}。

在条件~\ref{cond:7-2-18} 下，
\textbf{[AuBr95]} 的做法是取
\[
C_{\mathrm{dr}}(M,f)
=\bigoplus_i\Omega^j\Gamma(R_i),
\]
这里 $\Omega^j\Gamma(R_i)$
表示 $R_i$ 上所有光滑 $j$-形式，
然后定义
\begin{equation}
d_+(u)=du
+\sum_{j\ge i}
\bigl(\operatorname{ev}^{i,j}_{+\infty}\bigr)_!
\bigl(\operatorname{ev}^{i,j}_{-\infty}\bigr)^*(u).
\tag{7.2.19}\label{eq:7-2-19}
\end{equation}
这里
$\bigl(\operatorname{ev}^{i,j}_{-\infty}\bigr)^*$
是拉回，
$\bigl(\operatorname{ev}^{i,j}_{+\infty}\bigr)_!$
是沿纤维积分。
条件~\ref{cond:7-2-18}
保证沿纤维积分有定义，
并给出光滑形式。
以式~\eqref{eq:7-2-19}
代替式~\eqref{eq:7-2-17}，
便得到谱序列。

\begin{remark}{7.2.20}\label{rem:7-2-20}
近十年的技术发展使得无须假设
条件~\ref{cond:7-2-18}，
也可以使用德拉姆理论。
利用连续扰动族的方法
（见第~7.5 节），
可找到
$U_{i,j}\supset\mathcal M(R_i,R_j)$
及其上的光滑形式 $\omega_{i,j}$，
使
$\operatorname{ev}^{i,j}\colon
U_{i,j}\to R_i\times R_j$
为浸没，
$\omega_{i,j}$ 有紧支撑，
并使它成为流形链
$\mathcal M(R_i,R_j)$
所定义流的光滑化。
于是将~\eqref{eq:7-2-19}
改为
\[
d_+(u)=du+
\sum_{j\ge i}
\bigl(\operatorname{ev}^{i,j}_{+\infty}\bigr)_!
\left(
\bigl(\operatorname{ev}^{i,j}_{-\infty}\bigr)^*(u)
\wedge\omega_{i,j}
\right).
\]
对 $\omega_{i,j}$
施加适当的相容条件，
并配置适当符号，
即可使 $d_+\circ d_+=0$。
这里不再沿此方向展开，
细节留待读者完成。
\end{remark}

现在离开博特--莫尔斯理论的模型情形，
转向不再满足类似
式~\eqref{eq:7-2-11}
性质的更一般情形。
博特--莫尔斯型诺维科夫同调\liftIndex{I346}
就是一例，下面略作讨论。

设 $u$ 是有限维闭流形 $M$
上的闭 $1$-形式。
假定对 $M$ 的每个可缩子集 $U$，
都有 $u=df$，
其中 $f$ 是博特--莫尔斯函数。
定义 $u$ 的零点集
\[
\operatorname{Zero}(u)
=\{p\in M\mid u(p)=0\}
=\bigcup_iR_i.
\]
用向量场 $u^*$ 代替
$\operatorname{grad}f$，
同样定义
$\mathcal M(R_i,R_j)$；
这里 $u^*$ 由
\[
g(u^*,v)=u(v)
\]
刻画。
仍可在 $\operatorname{Zero}(u)$ 外
扰动 $u$，
使 $\mathcal M(R_i,R_j)$
成为带角流形。

与前例的关键区别是，
此时~\eqref{eq:7-2-11}
不再成立，
因而无法自然地给各 $R_i$ 编号，
以便沿用引理~\ref{lem:7-2-12}
证明中的归纳法。
这迫使我们使用更复杂的论证。
先取 $S(R_i;\mathbb Z)$
的有限生成子复形
$C_1(R_i;\mathbb Z)$，
使 (7.2.13.1) 成立；
然后扰动 $\mathcal M(R_i,R_j)$，
使它与这些链满足类似
(7.2.13.2) 的横截性条件。
但 (7.2.13.3) 随之未必成立，
于是扩大为
\[
C_{(2)}(R_i;\mathbb Z)
=C_1(R_i;\mathbb Z)
\oplus C_2(R_i;\mathbb Z),
\]
使 (7.2.13.1) 仍成立。
再对 $\mathcal M(R_i,R_j)$
另作一次扰动，
以取得这个新扰动与
$C_2(R_i;\mathbb Z)$
各元素间的 (7.2.13.2) 横截性。
然而 (7.2.13.3) 又未必成立，
故须加入第三项
$C_3(R_i;\mathbb Z)$，
如此继续。
最终得到下述命题。

\begin{proposition}{7.2.21}\label{prop:7-2-21}
对任意 $G$，
存在 $g=1,2,\ldots,G$
的奇异单形集
$\mathcal X_g(R_i)$，
以及依赖于 $g$ 的
$\mathcal M(R_i,R_j)$
的扰动
$\mathcal M(R_i,R_j)^g$，
满足下列性质。
记 $C_g(R_i;\mathbb Z)$ 为
$S(R_i;\mathbb Z)$
中由
$\bigcup_{h\le g}\mathcal X_h(R_i)$
生成的子复形。

\noindent\textup{(7.2.22.1)}\quad
自然包含
$C_g(R_i;\mathbb Z)
\to S(R_i;\mathbb Z)$
是链同伦等价。

\noindent\textup{(7.2.22.2)}\quad
$\mathcal X_g(R_i)$
的每个元素 $(P,f)$
都与
\[
\operatorname{ev}^{i,j}_{-\infty}
\colon S_m\mathcal M(R_i,R_j)^g
\longrightarrow R_i
\]
横截；这里 $j$ 任意，
$m=0,1,2,\ldots$。

\noindent\textup{(7.2.22.3)}\quad
对 $\mathcal X_g(R_i)$
的每个元素 $(P,f)$，
纤维积所得的链
\[
\bigl(
P_f\mathop{\times}_{\operatorname{ev}^{i,j}_{-\infty}}
S_m\mathcal M(R_i,R_j)^g,
\operatorname{ev}^{i,j}_{+\infty}
\bigr)
\in C_{g+1}(R_j;\mathbb Z)
\]
对任意 $j$ 及
$m=0,1,2,\ldots$
都可分解为
$\mathcal X_{g+1}(R_j)$
中的元素。
\translatornote{原书只定义了
$g=1,\ldots,G$ 的
$\mathcal X_g$、$C_g$，
但本条件在 $g=G$ 时
仍引用 $\mathcal X_{G+1}$、$C_{G+1}$。
是把条件限于 $g<G$，
还是多构造一级，
源书未作裁决；此处保留边界缺口。
参见疑误 A-0433。}

\noindent\textup{(7.2.22.4)}\quad
由式~\eqref{eq:7-2-17}
定义的算子 $\partial^+$
满足 $\partial^+\circ\partial^+=0$。
\end{proposition}

这里略去命题~\ref{prop:7-2-21}
的详细证明，
因为它与后文命题 7.2.35 的证明相似，
而且较为简单；
后一个命题将给出详细证明。
至于为何必须固定 $G$，
并在归纳的第 $G$ 级停止构造，
下一小节将作解释。
利用命题~\ref{prop:7-2-21}，
可以为诺维科夫同调构造
类似~\eqref{eq:7-2-16}
的谱序列。
还应指出，
结合 \textbf{[FuOn99II]} 的方法，
这一论证可以推广，
以直接实现皮乌尼欣
关于量子上同调与弗洛尔上同调等价性
的原始证明方案 \textbf{[Piu94]}。

\begin{remark}{7.2.23}\label{rem:7-2-23}
上文已经提到，
连续扰动族的方法
可在 D. Austin 与 P. Braam 的做法中
去掉条件~\ref{cond:7-2-18}。
只要使用 $\mathbb R$ 系数，
此法也给出处理本小节一般困难的
另一途径：
在德拉姆理论框架中完成构造。
连续扰动族的用法见第~7.5 节。
本书除第~7.5 节外
不采用这种德拉姆方法，
因为若干应用需要
$\mathbb Z$ 或 $\mathbb Z_2$ 系数。
连续扰动族的方法
将在别处详述
（见 \textbf{[Fuk07I]} 第 12 节及
\textbf{[FOOO08II]} 第 12 节）。
\end{remark}
\liftIndexRangeEnd{I384}
